\lecture{16}{28 April.}{}
Once we can interchange an integral with a sum of two functions, we can handle fininte sums by induction. More suprisingly, we can handle infinite sums of non-negative functions as well:

\begin{corollary} \label{cl: int and countable sum can interchange}
    If \(\Omega \subseteq \mathbb{R} ^n\) is measurable and \(g_1, g_2, \dots \) are a sequence of non-negative measurable functions from \(\Omega\) to \([0, \infty]\), then 
    \[
        \int _\Omega \sum_{n=1}^{\infty} g_n = \sum_{n=1}^{\infty} \int _\Omega g_n.  
    \]    
\end{corollary}

\begin{proof}
    For each integer \(N \ge 1\), define \(G_N \coloneqq \sum_{n=1}^N g_n \). Since \(g_n\) is measurable and non-negative, it follows that each \(G_N\) is also measurable and non-negative. Moreover, the sequence \((G_N)\) is increasing pointwise, because 
    \[
        G_{N+1}(x) = G_N(x) + g_{N+1}(x) \ge G_N(x) \quad \forall x \in \Omega.
    \]     
    Hence, we may apply the Lebesgue monotone convergence theorem to the sequence \(\left( G_N \right) \). We obtain 
    \[
        \int _\Omega \sup _N G_N = \sup _N \int _\Omega G_N.
    \] 
    Now since for each \(x \in \Omega\), \(G_1(x) \le G_2(x) \le \dots \), so \(\sup _N G_N(x) = \lim_{N \to \infty} G_N(x) = \sum_{n=1}^{\infty} g_n(x)  \). This gives
    \[
        \int _\Omega \sup _N G_N = \int _\Omega \sum_{n=1}^{\infty} g_n. 
    \]  
    Now since for each \(N \in \mathbb{N}\),
    \[
        \int _\Omega G_N = \int _\Omega \sum_{n=1}^N g_n = \sum_{n=1}^N \int _\Omega g_n.  
    \] 
    Hence,
    \[
        \sup _N \int _\Omega G_N = \sup _N \sum_{n=1}^N \int _\Omega g_n. 
    \]
    Note that the sequence \(\left( \sum_{n=1}^N \int _\Omega g_n  \right)_{N=1}^{\infty} \) is increasing, so 
    \[
        \sup _N \sum_{n=1}^N \int _\Omega g_n = \sum_{n=1}^{\infty} \int _\Omega g_n.  
    \] 
    Thus,
    \[
        \int _\Omega \sum_{n=1}^{\infty} g_n = \int _\Omega \sup _N G_N = \sup _N \int _\Omega G_N = \sum_{n=1}^{\infty} \int _\Omega g_n. 
    \]
\end{proof}

\begin{remark}
    We do not need to assume anything about convergence of above sums; it may well happen that both sides equal \(+\infty\). However, we do need to assume non-negativity.  
\end{remark}

One might similarly ask whether we can interchange limits and integrals; i.e. whether 
\[
    \int _\Omega \lim_{n \to \infty} f_n = \lim_{n \to \infty} \int _\Omega f_n  
\]
holds in general. Unfortunately this is not true: For \(n = 1, 2, 3, \dots \) let \(f_n: \mathbb{R} \to \mathbb{R} \) be \(f_n = \chi _{[n, n+1)}\). Then, \(f_n(x) \to 0\) for every \(x\), but \(\int _\mathbb{R} f_n = 1\) for every \(n\) since \(f_n\) is a simple function, and thus 
\[
    \lim_{n \to \infty} \int _\mathbb{R} f_n = 1 \neq 0. 
\]        

However, the following lemma of Fatou shows that the reverse cannot happen: the limiting function cannot have larger integral than the (liminf of the) original integrals. 

\begin{lemma}[Fatou's lemma]
    Let \(\Omega\) be a measurable subset of \(\mathbb{R} ^n\), and let \(f_1, f_2, \dots : \Omega \to [0, \infty]\) be a sequence of non-negative measurable functions. Then 
    \[
        \int _\Omega \liminf_{n \to \infty} f_n \le \liminf_{n \to \infty} \int _\Omega f_n.  
    \]   
\end{lemma}

\begin{proof}
    For each \(N \in \mathbb{N} \), define 
    \[
        g_N(x) \coloneqq \inf_{n \ge N} f_n(x) \quad \forall x \in \Omega.
    \] 
    Then \(g_N\) is non-negative and measurable. Moreover, \(g_1(x) \le g_2(x) \le \dots \) for all \(x \in \Omega\). By the definition of liminf, 
    \[
        \liminf_{n \to \infty} f_n(x) = \sup _{N \ge 1} g_N(x). 
    \]  
    Since \((g_N(x))\) is increasing, this means that \(g_N(x) \uparrow \liminf_{n \to \infty} f_n(x) \) for all \(x \in \Omega\). Note that we can apply Lebesgue monotone convergence theorem since \((g_N)\) is increasing:
    \[
        \int _\Omega \liminf_{n \to \infty} f_n(x) = \int _\Omega \sup _{N \ge 1} g_N = \sup _{N \ge 1} \int _\Omega g_N.
    \]
    On the other hand, for every \(N\) and every \(n \ge N\), we have \(g_N \le f_n\). Hence,
    \[
        \int _\Omega g_N \le \int _\Omega f_n \quad \forall n \ge N.
    \]  
    Taking infinimum over all \(n \ge N\), we get \(\int _\Omega g_N \le \inf_{n \ge N} \int _\Omega f_n\). Now take the supremum over \(N\) to obtain 
    \[
        \sup _{N \ge 1} \int _\Omega g_N \le \sup _{N \ge 1} \inf_{n \ge N} \int _\Omega f_n.
    \]   
    Combining previous result, we have 
    \[
       \int _\Omega \liminf_{n \to \infty} f_n(x) = \sup _{N \ge 1} \int _\Omega g_N \le \sup _{N \ge 1} \inf _{n \ge N} \int_\Omega f_n = \liminf_{n \to \infty} f_n.  
    \]

\end{proof}

Note that we are allowing our functions to take the value \(+ \infty\) at some points. It is even possible for a function to take the value \(+ \infty\) but still have a finite integral; for instance, if \(E\) is a set of measure zero, and \(f: \Omega \to \mathbb{R} \) is equal to \(+\infty \) on \(E\) but equal \(0\) everywhere else, then \(\int _\Omega f = 0\) by (a) of \autoref{prop: measurable func integral property}. However, if the integral is finite, the function must be finite almost everywhere:

\begin{lemma} \label{lm: finite integral implies infty set has measure 0}
    Let \(\Omega \) be a measurable subset of \(\mathbb{R} ^n\), and let \(f: \Omega \to [0, \infty]\) be a non-negative measurable function s.t. \(\int _\Omega f\) is finite. Then \(f\) is finite almost everywhere, i.e. the set \(\left\{ x \in \Omega : f(x) = + \infty  \right\} \) has measure zero.      
\end{lemma}

\begin{proof}
    For each integer \(k \ge 1\), define \(A_k \coloneqq \left\{ x \in \Omega: f(x) \ge k \right\} \). Since \(f\) is measurable and \([k, \infty]\) is a Borel subset of \([0, \infty]\), each set \(A_k\) is measurable. 

    On the set \(A_k\), we have \(f(x) \ge k\), while outside \(A_k\) the function \(k \chi _{A_k}\) vanishes. Thus, for every \(x \in \Omega\), \(k \chi _{A_k}(x) \le f(X)\). Because both \(k \chi _{A_k}\) and \(f\) are non-negative measurable functions, the monotonicity property gives \(\int _\Omega k \chi _{A_k} \le \int _\Omega f\). Now 
    \[
        \int _\Omega k \chi _{A_k} = k \int _\Omega \chi _{A_k} = k m(A_k),
    \]                
    so we obtain \(k m(A_k) \le \int _\Omega f\). Since \(\int _\Omega f < \infty\), it follows that \(m(A_k) \le \frac{1}{k} \int _\Omega f\), and thus 
    \[
        m(A_k) \to 0 \quad \text{as } k \to \infty. 
    \]   
    Note that 
    \[
        \left\{ x \in \Omega: f(x) = +\infty \right\} = \bigcap_{k=1}^{\infty} A_k.  
    \]
    Also, \(A_1 \supseteq A_2 \supseteq A_3 \supseteq \dots \) by definition, and by continuity of measure (HW5 P2),
    \[
        m \left( \bigcap_{k=1}^{\infty} A_k  \right) = \lim_{k \to \infty} m(A_k).  
    \]
    Since we have shown \(\lim_{k \to \infty} m(A_k) = 0 \), so we have 
    \[
        m\left( \left\{ x \in \Omega : f(x) = +\infty \right\}  \right) = m \left( \bigcap_{k=1}^{\infty} A_k  \right) = 0.  
    \] 
\end{proof}

\begin{lemma}[Borel-Cantelli lemma]
    Let \(\Omega _1, \Omega _2, \dots \) be measurable subsets of \(\mathbb{R} ^n\) such that
    \[
        \sum_{n=1}^{\infty} m(\Omega _n) < \infty.
    \]
    Then the set 
    \[
        \left\{ x \in \mathbb{R} ^n: x \in \Omega _k \text{ for infinitely many } k  \right\} 
    \]
    has measure zero. Equivalently, for almost every \(x \in \mathbb{R} ^n\), there exists \(N = N(x)\) s.t. \(x \notin \Omega _k\) for all \(k \ge N\).    
\end{lemma}

\begin{proof}
    Define 
    \[
        G(x) \coloneqq \sum_{n=1}^{\infty} \chi _{\Omega _n} (x), \quad \forall x \in \mathbb{R} ^n,
    \]
    where the value \(+ \infty\) is allowed. 

    For each \(n\), \(\chi _{\Omega _n}\) is measurable and non-negative. Hence, the partial sums 
    \[
        G_N(x) \coloneqq \sum_{n=1}^N \chi _{\Omega _n} (x) 
    \]
    are non-negative measurable functions. Since \(G_N(x) \uparrow G(x)\) for every \(x \in \mathbb{R} ^n\). It follows that \(G = \sup _N G_N\) is also measurable and non-negative. By \autoref{cl: int and countable sum can interchange}, we have 
    \[
        \int _{\mathbb{R} ^n} G = \int _{\mathbb{R} ^n} \sum_{n=1}^{\infty} \chi _{\Omega _n} = \sum_{n=1}^{\infty} \int _{\mathbb{R} ^n} \chi _{\Omega _n}.  
    \]   
    Since \(\int _{\mathbb{R} ^n} \chi _{\Omega _n} = m(\Omega _n)\), this becomes
    \[
        \int _{\mathbb{R} ^n} G = \sum_{n=1}^{\infty} m(\Omega _n) < \infty. 
    \] 
    By \autoref{lm: finite integral implies infty set has measure 0}, 
    \[
        m \left( \left\{ x \in \mathbb{R} ^n: G(x) = +\infty \right\}  \right) = 0. 
    \] 
    Note that 
    \[
        G(x) = +\infty \iff \sum_{n=1}^{\infty} \chi _{\Omega _n}(x) = \infty \iff x \in \Omega _k \text{ for infinitely many } k,  
    \]
    so 
    \[
        \left\{ x \in \mathbb{R} ^n: G(x) = +\infty \right\} = \left\{ x \in \mathbb{R} ^n: x \in \Omega _k \text{ for infinitely many } k  \right\}, 
    \]
    and we're done.
\end{proof}

\section{Integration of Absolutely Integrable Functions}
We have now completed the theory of the Lebesgue integral for non-negative functions. We now consider how to integrate functions which can be both positive and negative. However, we wish to avoid the indefinite expression \(+\infty + (- \infty)\), so we restrict our attention to a subclass of measurable functions: the absolutely integrable functions.

\begin{definition}
    Let \(\Omega\) be a measurable subset of \(\mathbb{R} ^n\). A measurable function \(f: \Omega \to \mathbb{R} ^*\) is said to be absolutely integrable if the integral \(\int _\Omega \vert f \vert \) is finite.
\end{definition}

Of course, \(\vert f \vert \) is always non-negative, so this definition makes sense even if \(f\) changes sign. Absolutely integrable functions are also known as \(L^1(\Omega)\) functions. 

If \(f: \Omega \to \mathbb{R} ^*\) is a function, we define the positive part \(f^+: \Omega \to [0, \infty]\) and the negative part \(f^-: \Omega \to [0, \infty]\) by 
\[
    f^+ \coloneqq \max (f, 0), \quad f^- \coloneqq - \min (f, 0).
\]   
Equivalently, for each \(x \in \Omega\), 
\[
    f^+(x) = \max \left\{ f(x), 0 \right\}, \quad f^-(x) = \max \left\{ -f(x), 0 \right\}.  
\] 
From \autoref{cl: measurable func plus, minus, ... is measurable}, which can be extended to \(\mathbb{R} ^*\)-valued functions without difficulty, we know that \(f^+, f^-\) are measurable. Clearly \(f^+\) and \(f^-\) are non-negative.

We also have the pointwise identities
\[
    f = f^+ - f^-, \quad \vert f \vert = f^+ + f^-. 
\]

\begin{definition}[Lebesgue integral]
    Let \(f: \Omega \to \mathbb{R} ^*\) be an absolutely integrable function. We define the Lebesgue integral \(\int _\Omega f\) of \(f\) to be the quantity
    \[
        \int _\Omega f \coloneqq \int _\Omega f^+ - \int _\Omega f^-.
    \]   
\end{definition}

\begin{remark}
    Note that since \(f\) is absolutely integrable, \(\int _\Omega f^+\) and \(\int _\Omega f^-\) are less than or equal to \(\int _\Omega \vert f \vert \) and hence are finite. Thus, \(\int _\Omega f\) is always finite; we are never encountering the indeterminate form \(+ \infty - (+ \infty)\).      
\end{remark}

\begin{remark}
    We have the useful triangle inequality
    \[
        \left\vert \int _\Omega f \right\vert \le \int _\Omega \vert f \vert.
    \]
    This is because
    \[
        \left\vert \int _\Omega f \right\vert = \left\vert \int _\Omega f^+ - \int _\Omega f^- \right\vert \le \left\vert \int _\Omega f^+ \right\vert + \left\vert \int _\Omega f^- \right\vert = \int _\Omega f^+ + \int _\Omega f^- = \int _\Omega f^+ + f^- = \int _\Omega \vert f \vert.     
    \]
\end{remark}